% !TeX root = ../../Book1.tex
% 纲要标题：### 17.1 有限域的存在与唯一性
% 纲要位置：计划纲要.md，第 587 行。

\section{有限域的存在与唯一性}
\label{b1:sec:1701}

% 纲要内容（仅作写作计划，尚非教材正文）：
%
% * 素数幂阶的必要性
% * 分裂域构造有限域；由 Frobenius 迭代说明根集本身闭成子域
% * 给定阶的有限域在同构意义下唯一
%

有限域不能有任意多个元素。例如六个元素的交换环 \(\mathbb Z/6\mathbb Z\) 含有零因子，因而不是域。限制域的大小的首先是素域和向量空间维数；反过来，给定一个允许的大小，分裂域把所需的全部元素一次构造出来。

\subsection{有限域阶的素数幂条件}
\label{b1:subsec:1701:01}

设 \(F\) 是有限域。它的特征不可能为零，否则整数的像给出无限多个不同元素。其特征是某个素数 \(p\)，素域为 \(\mathbb F_p\)。从线性无关元逐次扩充，若它们尚不生成 \(F\)，就选一个不在其线性生成空间中的元素加入。每次选入的元素都不同，而 \(F\) 是有限集合，所以这个过程必定停止，所得有限线性无关集又生成 \(F\)。故 \(n=[F:\mathbb F_p]\) 是正整数。取一组基 \(e_1,\ldots,e_n\)，每个元素唯一写成
\[
a_1e_1+\cdots+a_ne_n,\qquad a_i\in\mathbb F_p.
\]
基表示的存在性与唯一性说明，\(F\) 的元素与坐标向量 \((a_1,\ldots,a_n)\in\mathbb F_p^n\) 一一对应。每个坐标有 \(p\) 种选择，所以 \(|F|=p^n\)。特别地，六元域、十二元域均不存在；八元域与九元域则尚须构造，不能由必要条件直接断言存在。

\ProseParagraphBreak
更一般地，若 \(K\subseteq F\) 是有限域，\(|K|=q\)，\([F:K]=m\)，同样的基计数给出 \(|F|=q^m\)。这里 \(q\) 本身已经是素数幂。

\subsection{有限域的分裂域构造}
\label{b1:subsec:1701:02}

\begin{theorem}\label{b1:thm:finite-field-existence}
对每个素数 \(p\) 和每个正整数 \(n\)，存在一个恰有 \(p^n\) 个元素的域。它可取为 \(x^{p^n}-x\) 在 \(\mathbb F_p\) 上的分裂域。
\end{theorem}
\begin{proof}
记 \(q=p^n\)，取 \(x^q-x\) 的一个分裂域 \(L\)，并令 \(F\) 为它在 \(L\) 中的根集。需要证明这些根在域运算下仍然是根，才能把根集看成一个域。特征 \(p\) 下二项式展开给出 \((a+b)^p=a^p+b^p\)，再次取 \(p\) 次幂便有
\[
(a+b)^{p^2}=(a^p+b^p)^p=a^{p^2}+b^{p^2}.
\]
这样迭代 \(n\) 次，得到 \((a+b)^q=a^q+b^q\)。取幂也保持乘法和单位，所以 \(a\mapsto a^q\) 是环同态；它因而保持加法逆元，给出 \((-a)^q=-a^q\)。因此若 \(a^q=a,b^q=b\)，则
\[
(a+b)^q=a+b,\qquad (ab)^q=ab,\qquad (-a)^q=-a.
\]
零和一属于 \(F\)；若 \(a\ne0\)，则 \((a^{-1})^q=(a^q)^{-1}=a^{-1}\)。所以 \(F\) 对域的运算封闭，是 \(L\) 的子域，并且包含素域 \(\mathbb F_p\)。它已经含有全部根，便含有这些根在 \(\mathbb F_p\) 上生成的域，而这个域就是分裂域 \(L\)；结合 \(F\subseteq L\)，得到 \(L=F\)。另一方面，\(q=p^n\) 在特征 \(p\) 中为零，所以 \(x^q-x\) 的导数为 \(-1\)，没有重根；它在 \(L\) 中分裂为 \(q\) 个互异的一次因子，所以 \(|F|=q\)。
\end{proof}
例如 \(x^2+x+1\) 在 \(\mathbb F_2\) 上没有根，因而不可约。在 \(\mathbb F_2[x]\) 中有
\[
x^4-x=x(x+1)(x^2+x+1).
\]
商域 \(\mathbb F_2[x]/(x^2+x+1)\) 有四个元素。记 \(x\) 的剩余类为 \(\alpha\)，则不可约二次因子的两个根为 \(\alpha,\alpha+1\)，加上零和一恰为全部四个元素。这个商域由 \(\alpha\) 生成，也是 \(x^4-x\) 的分裂域，正是上述构造在 \(q=4\) 时的一种具体表示。

\subsection{给定阶有限域的唯一性}
\label{b1:subsec:1701:03}

\begin{theorem}\label{b1:thm:finite-field-unique}
设 \(p\) 为素数，\(n\geq1\) 为整数，并令 \(q=p^n\)。任意两个具有 \(q\) 个元素的域，都在其素域 \(\mathbb F_p\) 上同构。
\end{theorem}
\begin{proof}
若 \(|F|=q=p^n\)，由已证的素数幂条件及整数素因子分解的唯一性，\(F\) 的素域是 \(\mathbb F_p\)。其乘法群 \(F^\times\) 有 \(q-1\) 个元素。由\cref{b1:thm:lagrange}，每个非零 \(a\in F\) 都满足 \(a^{q-1}=1\)，故 \(F\) 的全部 \(q\) 个元素都是 \(x^q-x\) 的根。次数为 \(q\) 的多项式不可能有更多根，因此这些元素就是它在任意扩张域中的全部根。根集既然等于整个 \(F\)，由根在素域上生成的域也等于 \(F\)，所以 \(F\) 是该多项式在 \(\mathbb F_p\) 上的分裂域。\cref{b1:prop:splitting-field-unique}保证同一个多项式的两个分裂域在保持基域的同构意义下唯一，故两个这样的域同构。
\end{proof}
以后用 \(\mathbb F_q\) 表示任一具有 \(q\) 个元素的域。
\symidx{\(\mathbb F_q\)：具有 \(q\) 个元素的有限域}
这个记号指定的是同构类型，并不指定某一组元素名称，也不指定唯一同构。例如四元域中的两个非素域元素可以交换。未选定共同的环境域时，两个这样的符号也不自带集合包含关系。需要比较多个有限域的包含关系时，本书先固定 \(\overline{\mathbb F}_p\)，再把 \(\mathbb F_{p^n}\) 取为其中 \(x^{p^n}-x\) 的根集；此时它才是环境域中的确定子域，后面的包含、交与合成均在这个意义下使用。
